\documentclass[
  man,
  floatsintext,
  longtable,
  nolmodern,
  notxfonts,
  notimes,
  colorlinks=true,linkcolor=blue,citecolor=blue,urlcolor=blue]{apa7}

\usepackage{amsmath}
\usepackage{amssymb}




\RequirePackage{longtable}
\RequirePackage{threeparttablex}

\makeatletter
\renewcommand{\paragraph}{\@startsection{paragraph}{4}{\parindent}%
	{0\baselineskip \@plus 0.2ex \@minus 0.2ex}%
	{-.5em}%
	{\normalfont\normalsize\bfseries\typesectitle}}

\renewcommand{\subparagraph}[1]{\@startsection{subparagraph}{5}{0.5em}%
	{0\baselineskip \@plus 0.2ex \@minus 0.2ex}%
	{-\z@\relax}%
	{\normalfont\normalsize\bfseries\itshape\hspace{\parindent}{#1}\textit{\addperi}}{\relax}}
\makeatother




\usepackage{longtable, booktabs, multirow, multicol, colortbl, hhline, caption, array, float, xpatch}
\usepackage{subcaption}


\renewcommand\thesubfigure{\Alph{subfigure}}
\setcounter{topnumber}{2}
\setcounter{bottomnumber}{2}
\setcounter{totalnumber}{4}
\renewcommand{\topfraction}{0.85}
\renewcommand{\bottomfraction}{0.85}
\renewcommand{\textfraction}{0.15}
\renewcommand{\floatpagefraction}{0.7}

\usepackage{tcolorbox}
\tcbuselibrary{listings,theorems, breakable, skins}
\usepackage{fontawesome5}

\definecolor{quarto-callout-color}{HTML}{909090}
\definecolor{quarto-callout-note-color}{HTML}{0758E5}
\definecolor{quarto-callout-important-color}{HTML}{CC1914}
\definecolor{quarto-callout-warning-color}{HTML}{EB9113}
\definecolor{quarto-callout-tip-color}{HTML}{00A047}
\definecolor{quarto-callout-caution-color}{HTML}{FC5300}
\definecolor{quarto-callout-color-frame}{HTML}{ACACAC}
\definecolor{quarto-callout-note-color-frame}{HTML}{4582EC}
\definecolor{quarto-callout-important-color-frame}{HTML}{D9534F}
\definecolor{quarto-callout-warning-color-frame}{HTML}{F0AD4E}
\definecolor{quarto-callout-tip-color-frame}{HTML}{02B875}
\definecolor{quarto-callout-caution-color-frame}{HTML}{FD7E14}

%\newlength\Oldarrayrulewidth
%\newlength\Oldtabcolsep


\usepackage{hyperref}



\usepackage{color}
\usepackage{fancyvrb}
\newcommand{\VerbBar}{|}
\newcommand{\VERB}{\Verb[commandchars=\\\{\}]}
\DefineVerbatimEnvironment{Highlighting}{Verbatim}{commandchars=\\\{\}}
% Add ',fontsize=\small' for more characters per line
\usepackage{framed}
\definecolor{shadecolor}{RGB}{241,243,245}
\newenvironment{Shaded}{\begin{snugshade}}{\end{snugshade}}
\newcommand{\AlertTok}[1]{\textcolor[rgb]{0.68,0.00,0.00}{#1}}
\newcommand{\AnnotationTok}[1]{\textcolor[rgb]{0.37,0.37,0.37}{#1}}
\newcommand{\AttributeTok}[1]{\textcolor[rgb]{0.40,0.45,0.13}{#1}}
\newcommand{\BaseNTok}[1]{\textcolor[rgb]{0.68,0.00,0.00}{#1}}
\newcommand{\BuiltInTok}[1]{\textcolor[rgb]{0.00,0.23,0.31}{#1}}
\newcommand{\CharTok}[1]{\textcolor[rgb]{0.13,0.47,0.30}{#1}}
\newcommand{\CommentTok}[1]{\textcolor[rgb]{0.37,0.37,0.37}{#1}}
\newcommand{\CommentVarTok}[1]{\textcolor[rgb]{0.37,0.37,0.37}{\textit{#1}}}
\newcommand{\ConstantTok}[1]{\textcolor[rgb]{0.56,0.35,0.01}{#1}}
\newcommand{\ControlFlowTok}[1]{\textcolor[rgb]{0.00,0.23,0.31}{\textbf{#1}}}
\newcommand{\DataTypeTok}[1]{\textcolor[rgb]{0.68,0.00,0.00}{#1}}
\newcommand{\DecValTok}[1]{\textcolor[rgb]{0.68,0.00,0.00}{#1}}
\newcommand{\DocumentationTok}[1]{\textcolor[rgb]{0.37,0.37,0.37}{\textit{#1}}}
\newcommand{\ErrorTok}[1]{\textcolor[rgb]{0.68,0.00,0.00}{#1}}
\newcommand{\ExtensionTok}[1]{\textcolor[rgb]{0.00,0.23,0.31}{#1}}
\newcommand{\FloatTok}[1]{\textcolor[rgb]{0.68,0.00,0.00}{#1}}
\newcommand{\FunctionTok}[1]{\textcolor[rgb]{0.28,0.35,0.67}{#1}}
\newcommand{\ImportTok}[1]{\textcolor[rgb]{0.00,0.46,0.62}{#1}}
\newcommand{\InformationTok}[1]{\textcolor[rgb]{0.37,0.37,0.37}{#1}}
\newcommand{\KeywordTok}[1]{\textcolor[rgb]{0.00,0.23,0.31}{\textbf{#1}}}
\newcommand{\NormalTok}[1]{\textcolor[rgb]{0.00,0.23,0.31}{#1}}
\newcommand{\OperatorTok}[1]{\textcolor[rgb]{0.37,0.37,0.37}{#1}}
\newcommand{\OtherTok}[1]{\textcolor[rgb]{0.00,0.23,0.31}{#1}}
\newcommand{\PreprocessorTok}[1]{\textcolor[rgb]{0.68,0.00,0.00}{#1}}
\newcommand{\RegionMarkerTok}[1]{\textcolor[rgb]{0.00,0.23,0.31}{#1}}
\newcommand{\SpecialCharTok}[1]{\textcolor[rgb]{0.37,0.37,0.37}{#1}}
\newcommand{\SpecialStringTok}[1]{\textcolor[rgb]{0.13,0.47,0.30}{#1}}
\newcommand{\StringTok}[1]{\textcolor[rgb]{0.13,0.47,0.30}{#1}}
\newcommand{\VariableTok}[1]{\textcolor[rgb]{0.07,0.07,0.07}{#1}}
\newcommand{\VerbatimStringTok}[1]{\textcolor[rgb]{0.13,0.47,0.30}{#1}}
\newcommand{\WarningTok}[1]{\textcolor[rgb]{0.37,0.37,0.37}{\textit{#1}}}

\providecommand{\tightlist}{%
  \setlength{\itemsep}{0pt}\setlength{\parskip}{0pt}}
\usepackage{longtable,booktabs,array}
\usepackage{calc} % for calculating minipage widths
% Correct order of tables after \paragraph or \subparagraph
\usepackage{etoolbox}
\makeatletter
\patchcmd\longtable{\par}{\if@noskipsec\mbox{}\fi\par}{}{}
\makeatother
% Allow footnotes in longtable head/foot
\IfFileExists{footnotehyper.sty}{\usepackage{footnotehyper}}{\usepackage{footnote}}
\makesavenoteenv{longtable}

\usepackage{graphicx}
\makeatletter
\newsavebox\pandoc@box
\newcommand*\pandocbounded[1]{% scales image to fit in text height/width
  \sbox\pandoc@box{#1}%
  \Gscale@div\@tempa{\textheight}{\dimexpr\ht\pandoc@box+\dp\pandoc@box\relax}%
  \Gscale@div\@tempb{\linewidth}{\wd\pandoc@box}%
  \ifdim\@tempb\p@<\@tempa\p@\let\@tempa\@tempb\fi% select the smaller of both
  \ifdim\@tempa\p@<\p@\scalebox{\@tempa}{\usebox\pandoc@box}%
  \else\usebox{\pandoc@box}%
  \fi%
}
% Set default figure placement to htbp
\def\fps@figure{htbp}
\makeatother


% definitions for citeproc citations
\NewDocumentCommand\citeproctext{}{}
\NewDocumentCommand\citeproc{mm}{%
  \begingroup\def\citeproctext{#2}\cite{#1}\endgroup}
\makeatletter
 % allow citations to break across lines
 \let\@cite@ofmt\@firstofone
 % avoid brackets around text for \cite:
 \def\@biblabel#1{}
 \def\@cite#1#2{{#1\if@tempswa , #2\fi}}
\makeatother
\newlength{\cslhangindent}
\setlength{\cslhangindent}{1.5em}
\newlength{\csllabelwidth}
\setlength{\csllabelwidth}{3em}
\newenvironment{CSLReferences}[2] % #1 hanging-indent, #2 entry-spacing
 {\begin{list}{}{%
  \setlength{\itemindent}{0pt}
  \setlength{\leftmargin}{0pt}
  \setlength{\parsep}{0pt}
  % turn on hanging indent if param 1 is 1
  \ifodd #1
   \setlength{\leftmargin}{\cslhangindent}
   \setlength{\itemindent}{-1\cslhangindent}
  \fi
  % set entry spacing
  \setlength{\itemsep}{#2\baselineskip}}}
 {\end{list}}
\usepackage{calc}
\newcommand{\CSLBlock}[1]{\hfill\break\parbox[t]{\linewidth}{\strut\ignorespaces#1\strut}}
\newcommand{\CSLLeftMargin}[1]{\parbox[t]{\csllabelwidth}{\strut#1\strut}}
\newcommand{\CSLRightInline}[1]{\parbox[t]{\linewidth - \csllabelwidth}{\strut#1\strut}}
\newcommand{\CSLIndent}[1]{\hspace{\cslhangindent}#1}





\usepackage{newtx}

\defaultfontfeatures{Scale=MatchLowercase}
\defaultfontfeatures[\rmfamily]{Ligatures=TeX,Scale=1}





\title{Supplement: Parameter Recovery for a Custom \(M^3\) --- Memory
Updating Example}


\shorttitle{Supplement: \(M^3\) Parameter Recovery}


\usepackage{etoolbox}









\authorsnames[{1},{2,3}]{Gidon T. Frischkorn,Chenyu Li}







\authorsaffiliations{
{Department of Psychology, University of Zurich},{Department of
Psychology, University of Chicago},{Institution for Mind and
Biology,, University of Chicago}}




\leftheader{Frischkorn and Li}



\abstract{This supplement accompanies the tutorial on fitting the memory
measurement model (\(M^3\)) with the bmm R package. It presents a
complete parameter recovery simulation for a custom \(M^3\) applied to a
memory updating task with five response categories, five estimated
parameters, and experimentally manipulated design variables. The
walkthrough covers model definition, data simulation with
\texttt{rm3()}, model fitting, and recovery evaluation at both the group
and individual level. Results across a design grid varying sample size
and trial count illustrate how data quantity affects recovery quality.
Complete R scripts are available in the companion repository. }

\keywords{working memory, measurement model, parameter
recovery, simulation, R package}

\authornote{\par{\addORCIDlink{Gidon T.
Frischkorn}{0000-0002-5055-9764}}\par{\addORCIDlink{Chenyu
Li}{0000-0001-8815-6189}} 

\par{       }
\par{Correspondence concerning this article should be addressed to Gidon
T.
Frischkorn, Email: \href{mailto:gidon.frischkorn@psychologie.uzh.ch}{gidon.frischkorn@psychologie.uzh.ch}}
}

\makeatletter
\let\endoldlt\endlongtable
\def\endlongtable{
\hline
\endoldlt
}
\makeatother

\urlstyle{same}



\makeatletter
\@ifpackageloaded{caption}{}{\usepackage{caption}}
\AtBeginDocument{%
\ifdefined\contentsname
  \renewcommand*\contentsname{Table of contents}
\else
  \newcommand\contentsname{Table of contents}
\fi
\ifdefined\listfigurename
  \renewcommand*\listfigurename{List of Figures}
\else
  \newcommand\listfigurename{List of Figures}
\fi
\ifdefined\listtablename
  \renewcommand*\listtablename{List of Tables}
\else
  \newcommand\listtablename{List of Tables}
\fi
\ifdefined\figurename
  \renewcommand*\figurename{Figure}
\else
  \newcommand\figurename{Figure}
\fi
\ifdefined\tablename
  \renewcommand*\tablename{Table}
\else
  \newcommand\tablename{Table}
\fi
}
\@ifpackageloaded{float}{}{\usepackage{float}}
\floatstyle{ruled}
\@ifundefined{c@chapter}{\newfloat{codelisting}{h}{lop}}{\newfloat{codelisting}{h}{lop}[chapter]}
\floatname{codelisting}{Listing}
\newcommand*\listoflistings{\listof{codelisting}{List of Listings}}
\makeatother
\makeatletter
\makeatother
\makeatletter
\@ifpackageloaded{caption}{}{\usepackage{caption}}
\@ifpackageloaded{subcaption}{}{\usepackage{subcaption}}
\makeatother

% From https://tex.stackexchange.com/a/645996/211326
%%% apa7 doesn't want to add appendix section titles in the toc
%%% let's make it do it
\makeatletter
\xpatchcmd{\appendix}
  {\par}
  {\addcontentsline{toc}{section}{\@currentlabelname}\par}
  {}{}
\makeatother

%% Disable longtable counter
%% https://tex.stackexchange.com/a/248395/211326

\usepackage{etoolbox}

\makeatletter
\patchcmd{\LT@caption}
  {\bgroup}
  {\bgroup\global\LTpatch@captiontrue}
  {}{}
\patchcmd{\longtable}
  {\par}
  {\par\global\LTpatch@captionfalse}
  {}{}
\apptocmd{\endlongtable}
  {\ifLTpatch@caption\else\addtocounter{table}{-1}\fi}
  {}{}
\newif\ifLTpatch@caption
\makeatother

\begin{document}

\maketitle




\setlength\LTleft{0pt}




\section{Overview}\label{overview}

This supplement demonstrates a complete parameter recovery simulation
for a custom \(M^3\) using a memory updating task as a worked example.
The walkthrough follows the simplified companion script
(\texttt{scripts/tutorial5\_parameter\_recovery\_updating\_simple.R});
design-grid results from multiple sample sizes and trial counts come
from \texttt{scripts/tutorial5\_parameter\_recovery\_updating.R}.

The parameter recovery simulation workflow consists of four steps:

\begin{enumerate}
\def\labelenumi{\arabic{enumi}.}
\tightlist
\item
  Define a custom \(M^3\) with user-specified activation formulas.
\item
  Simulate data with known parameter values using \texttt{rm3()}.
\item
  Fit the model to the simulated data.
\item
  Compare recovered estimates to the true generating values.
\end{enumerate}

\subsection{Why parameter recovery}\label{why-parameter-recovery}

Before applying a new or modified \(M^3\) to empirical data, the model
should be able to recover the parameters that generated the data. If
recovery fails under ideal conditions (known generating values, no model
misspecification), there are likely problems with the model
specification or the design providing the data to identify the
parameters. Recovery simulations are especially informative for custom
models where the number of parameters and the functional form of the
activation equations are chosen by the researcher
(\citeproc{ref-wilson2019ten}{Wilson \& Collins, 2019}). They are also
useful for determining required sample size and the number of trials per
condition during the design phase of a study.

\section{\texorpdfstring{Defining a Custom \(M^3\) for Memory
Updating}{Defining a Custom M\^{}3 for Memory Updating}}\label{defining-a-custom-m3-for-memory-updating}

For this simulation we use a memory updating task as a running example
to illustrate a common situation: the task produces response categories
and activation structures that do not match either the \texttt{ss} or
\texttt{cs} versions of the \(M^3\). Specifically, replaced (outdated)
items create additional response categories, and two experimental timing
variables modulate activations in ways that require custom formulas. The
task design mimics the updating task from Oberauer and Lewandowsky
(\citeproc{ref-oberauer2019simple}{2019}). Participants memorize items
at five positions, then positions are updated randomly. At test, one
position is probed and the participant selects from a recognition set
containing: the current target (1 option), other current items (4
options), the old item from the tested position (1 option), old items
from other positions (4 options), and not-presented lures (5 options).
This produces five response categories with the following activation
formulas based on the model proposed by Oberauer and Lewandowsky
(\citeproc{ref-oberauer2019simple}{2019}):

\begin{equation}\phantomsection\label{eq-act-updating}{
\begin{aligned}
\text{A}_{\text{correct}}  &= b + (1 + e \cdot t_e) \cdot (a + c) \\
\text{A}_{\text{other}}    &= b + (1 + e \cdot t_e) \cdot a \\
\text{A}_{\text{oldinpos}} &= b + \exp(-r \cdot t_r) \cdot d \cdot (1 + e \cdot t_e) \cdot (a + c) \\
\text{A}_{\text{otherold}} &= b + \exp(-r \cdot t_r) \cdot d \cdot (1 + e \cdot t_e) \cdot a \\
\text{A}_{\text{NPL}}      &= b
\end{aligned}
}\end{equation}

Beyond the familiar \(a\) (item memory) and \(c\) (context binding)
parameters, three new parameters appear: \(d\) captures the residual
activation of replaced (outdated) items that represents instant
\emph{deletion} of outdated items independent of the time manipulations,
\(e\) scales the benefit of extended encoding time \(t_e\), and \(r\)
governs how old-item activation is removed with removal time \(t_r\). To
understand the composite terms: \((1 + e \cdot t_e)\) scales all memory
activations as a function of encoding time --- when encoding is brief
the boost is minimal, when it is long activations are amplified. The
term \(\exp(-r \cdot t_r) \cdot d\) governs how much activation old
items retain: \(d\) sets the residual activation at the moment of
replacement, and \(r\) controls how quickly this residual is removed
over time. Note that \(t_e\) and \(t_r\) are experimentally manipulated
design variables (in our simulation each at 0.25, 0.5, and 1.0 s), not
estimated parameters, producing a \(3 \times 3\) within-subjects design
with 9 conditions.

\subsection{Model specification in
bmm}\label{model-specification-in-bmm}

The model definition specifies the response categories and option
counts. In addition, we specify link functions to match each parameter's
theoretical constraints: identity for \(a\) and \(c\) because
activations are unbounded; logit for \(d\) because it represents a
proportion of original activation retained after replacement (bounded
between 0 and 1); and log for \(e\) and \(r\) because encoding benefit
and decay rate must be positive:

\begin{Shaded}
\begin{Highlighting}[]
\NormalTok{m3\_model\_custom }\OtherTok{\textless{}{-}} \FunctionTok{m3}\NormalTok{(}
  \AttributeTok{resp\_cats =} \FunctionTok{c}\NormalTok{(}\StringTok{"correct"}\NormalTok{, }\StringTok{"other"}\NormalTok{, }\StringTok{"oldinpos"}\NormalTok{, }\StringTok{"otherold"}\NormalTok{, }\StringTok{"npl"}\NormalTok{),}
  \AttributeTok{num\_options =} \FunctionTok{c}\NormalTok{(}
    \StringTok{"n\_correct"} \OtherTok{=} \DecValTok{1}\NormalTok{, }\StringTok{"n\_other"} \OtherTok{=} \DecValTok{4}\NormalTok{, }\StringTok{"n\_oldinpos"} \OtherTok{=} \DecValTok{1}\NormalTok{,}
    \StringTok{"n\_otherold"} \OtherTok{=} \DecValTok{4}\NormalTok{, }\StringTok{"n\_npl"} \OtherTok{=} \DecValTok{5}
\NormalTok{  ),}
  \AttributeTok{version     =} \StringTok{"custom"}\NormalTok{,}
  \AttributeTok{choice\_rule =} \StringTok{"softmax"}\NormalTok{,}
  \AttributeTok{links =} \FunctionTok{list}\NormalTok{(}
    \AttributeTok{a =} \StringTok{"identity"}\NormalTok{, }\AttributeTok{c =} \StringTok{"identity"}\NormalTok{,}
    \AttributeTok{d =} \StringTok{"logit"}\NormalTok{, }\AttributeTok{e =} \StringTok{"log"}\NormalTok{, }\AttributeTok{r =} \StringTok{"log"}
\NormalTok{  )}
\NormalTok{)}
\end{Highlighting}
\end{Shaded}

The activation formulas are specified separately using \texttt{bmf()}.
For the simulation we do not need to add predictor formulas for the
parameters, as data generating parameters will be passed directly to
\texttt{rm3()}:

\begin{Shaded}
\begin{Highlighting}[]
\NormalTok{m3\_act\_funs }\OtherTok{\textless{}{-}} \FunctionTok{bmf}\NormalTok{(}
\NormalTok{  correct  }\SpecialCharTok{\textasciitilde{}}\NormalTok{ (}\DecValTok{1} \SpecialCharTok{+}\NormalTok{ e }\SpecialCharTok{*}\NormalTok{ te) }\SpecialCharTok{*}\NormalTok{ (a }\SpecialCharTok{+}\NormalTok{ c) }\SpecialCharTok{+}\NormalTok{ b,}
\NormalTok{  other    }\SpecialCharTok{\textasciitilde{}}\NormalTok{ (}\DecValTok{1} \SpecialCharTok{+}\NormalTok{ e }\SpecialCharTok{*}\NormalTok{ te) }\SpecialCharTok{*}\NormalTok{ a }\SpecialCharTok{+}\NormalTok{ b,}
\NormalTok{  oldinpos }\SpecialCharTok{\textasciitilde{}} \FunctionTok{exp}\NormalTok{(}\SpecialCharTok{{-}}\NormalTok{r }\SpecialCharTok{*}\NormalTok{ tr) }\SpecialCharTok{*}\NormalTok{ d }\SpecialCharTok{*}\NormalTok{ (}\DecValTok{1} \SpecialCharTok{+}\NormalTok{ e }\SpecialCharTok{*}\NormalTok{ te) }\SpecialCharTok{*}\NormalTok{ (a }\SpecialCharTok{+}\NormalTok{ c) }\SpecialCharTok{+}\NormalTok{ b,}
\NormalTok{  otherold }\SpecialCharTok{\textasciitilde{}} \FunctionTok{exp}\NormalTok{(}\SpecialCharTok{{-}}\NormalTok{r }\SpecialCharTok{*}\NormalTok{ tr) }\SpecialCharTok{*}\NormalTok{ d }\SpecialCharTok{*}\NormalTok{ (}\DecValTok{1} \SpecialCharTok{+}\NormalTok{ e }\SpecialCharTok{*}\NormalTok{ te) }\SpecialCharTok{*}\NormalTok{ a }\SpecialCharTok{+}\NormalTok{ b,}
\NormalTok{  npl      }\SpecialCharTok{\textasciitilde{}}\NormalTok{ b}
\NormalTok{)}
\end{Highlighting}
\end{Shaded}

Note that \texttt{te} and \texttt{tr} appear in the formulas but are not
estimated parameters --- they are data columns whose values vary across
conditions. When \texttt{bmm} encounters variables in the activation
formulas that are not listed as model parameters, it treats them as data
columns and requires them to be present in the dataset. This allows
custom models to incorporate any number of experimentally manipulated
design variables directly into the activation equations.

\section{Choosing Generating Parameter
Values}\label{choosing-generating-parameter-values}

For the recovery simulation, we need group-level means and
between-person standard deviations for each parameter. These are
specified on the link scale (the scale on which the sampler operates),
not the native scale, to respect the required bounds. The values were
chosen to approximate empirical estimates from Li et al.
(\citeproc{ref-li2025updating}{2025}) and to produce realistic
individual differences. Table~\ref{tbl-gen-pars} lists the generating
values. The native range column shows the approximate span of individual
differences (\(\pm 2\) SD around the mean, transformed to native scale
via the inverse link function).

\begin{table}

{\caption{{Generating parameter values for the recovery simulation.
Parameters are defined on the link scale; the native range is the
approximate \(\pm 2\) SD interval after applying the inverse link
function.}{\label{tbl-gen-pars}}}
\vspace{-20pt}}

\begin{longtable}[]{@{}
  >{\raggedright\arraybackslash}p{(\linewidth - 8\tabcolsep) * \real{0.1528}}
  >{\raggedright\arraybackslash}p{(\linewidth - 8\tabcolsep) * \real{0.1389}}
  >{\raggedright\arraybackslash}p{(\linewidth - 8\tabcolsep) * \real{0.1806}}
  >{\raggedright\arraybackslash}p{(\linewidth - 8\tabcolsep) * \real{0.1528}}
  >{\raggedright\arraybackslash}p{(\linewidth - 8\tabcolsep) * \real{0.3750}}@{}}
\toprule\noalign{}
\begin{minipage}[b]{\linewidth}\raggedright
Parameter
\end{minipage} & \begin{minipage}[b]{\linewidth}\raggedright
Link
\end{minipage} & \begin{minipage}[b]{\linewidth}\raggedright
Mean (link)
\end{minipage} & \begin{minipage}[b]{\linewidth}\raggedright
SD (link)
\end{minipage} & \begin{minipage}[b]{\linewidth}\raggedright
Native range (\(\pm 2\) SD)
\end{minipage} \\
\midrule\noalign{}
\endhead
\bottomrule\noalign{}
\endlastfoot
\(c\) & identity & 2.0 & 0.50 & 1.0 -- 3.0 \\
\(a\) & identity & 1.0 & 0.50 & 0.0 -- 2.0 \\
\(d\) & logit & 0.85 & 0.65 & 0.38 -- 0.89 \\
\(e\) & log & 0.00 & 0.35 & 0.50 -- 2.01 \\
\(r\) & log & \$-\$0.69 & 0.38 & 0.24 -- 1.04 \\
\end{longtable}

\end{table}

Individual parameters are drawn from
\(\mathcal{N}(\mu_{\text{link}}, \sigma_{\text{link}})\) independently
for each person and parameter. We simulated \(N = 50\) participants with
25 observations per condition (225 observations total across the
\(3 \times 3\) design). Because the task uses set size 5 and we assume
that all positions within each trial are tested, each trial yields 5
responses, so 25 observations correspond to 5 trials per condition. With
9 conditions this amounts to 45 trials in total per participant,
corresponding to roughly one hour of testing.

\section{\texorpdfstring{Simulating Data with
\texttt{rm3()}}{Simulating Data with rm3()}}\label{simulating-data-with-rm3}

The \texttt{bmm} package provides the random generation function
\texttt{rm3()} as the data-generating counterpart to the model-fitting
workflow. Given a set of parameter values on the native scale, the model
definition, and the activation formulas, it computes category
probabilities via the specified choice rule and draws random samples
from the multinomial distribution. Here is an example call for
simulating a single participant in one condition:

\begin{Shaded}
\begin{Highlighting}[]
\NormalTok{demo\_data }\OtherTok{\textless{}{-}} \FunctionTok{rm3}\NormalTok{(}
  \AttributeTok{n =} \DecValTok{1}\NormalTok{, }\AttributeTok{size =} \DecValTok{25}\NormalTok{,}
  \AttributeTok{pars =} \FunctionTok{c}\NormalTok{(}\AttributeTok{a =} \FloatTok{1.2}\NormalTok{, }\AttributeTok{c =} \FloatTok{2.3}\NormalTok{, }\AttributeTok{d =} \FloatTok{0.7}\NormalTok{, }\AttributeTok{e =} \FloatTok{1.1}\NormalTok{, }\AttributeTok{r =} \FloatTok{0.5}\NormalTok{,}
           \AttributeTok{b =} \DecValTok{0}\NormalTok{, }\AttributeTok{te =} \FloatTok{0.25}\NormalTok{, }\AttributeTok{tr =} \FloatTok{0.25}\NormalTok{),}
  \AttributeTok{m3\_model =}\NormalTok{ m3\_model\_custom,}
  \AttributeTok{act\_funs =}\NormalTok{ m3\_act\_funs}
\NormalTok{)}
\end{Highlighting}
\end{Shaded}

This returns a one-row data frame with response counts across the five
categories. All variables that appear in the activation formulas must be
included in the \texttt{pars} vector: estimated parameters (\(a\),
\(c\), \(d\), \(e\), \(r\)), fixed parameters (\(b = 0\) for the softmax
choice rule), and design variables (\(t_e\), \(t_r\)). To simulate a
full dataset, we apply \texttt{rm3()} over all 50 participants and 9
conditions (\(3\) encoding times \(\times\) \(3\) removal times),
producing 450 rows of simulated data. The companion script contains the
complete generation workflow and data inspection code.

\section{Model Fitting}\label{model-fitting}

We fit the simulated data with the same \texttt{bmm()} workflow shown in
the main tutorial. For custom models, the prediction formula serves
double duty: it repeats the activation equations (telling \texttt{bmm}
how parameters combine into category activations) and appends regression
formulas that specify how each parameter varies across participants
(here, an intercept plus a random intercept per participant):

\begin{Shaded}
\begin{Highlighting}[]
\CommentTok{\# Prediction formula: intercept + random intercept per participant}
\NormalTok{m3\_pred\_funs }\OtherTok{\textless{}{-}} \FunctionTok{bmf}\NormalTok{(}
\NormalTok{  correct  }\SpecialCharTok{\textasciitilde{}}\NormalTok{ (}\DecValTok{1} \SpecialCharTok{+}\NormalTok{ e }\SpecialCharTok{*}\NormalTok{ te) }\SpecialCharTok{*}\NormalTok{ (a }\SpecialCharTok{+}\NormalTok{ c) }\SpecialCharTok{+}\NormalTok{ b,}
\NormalTok{  other    }\SpecialCharTok{\textasciitilde{}}\NormalTok{ (}\DecValTok{1} \SpecialCharTok{+}\NormalTok{ e }\SpecialCharTok{*}\NormalTok{ te) }\SpecialCharTok{*}\NormalTok{ a }\SpecialCharTok{+}\NormalTok{ b,}
\NormalTok{  oldinpos }\SpecialCharTok{\textasciitilde{}} \FunctionTok{exp}\NormalTok{(}\SpecialCharTok{{-}}\NormalTok{r }\SpecialCharTok{*}\NormalTok{ tr) }\SpecialCharTok{*}\NormalTok{ d }\SpecialCharTok{*}\NormalTok{ (}\DecValTok{1} \SpecialCharTok{+}\NormalTok{ e }\SpecialCharTok{*}\NormalTok{ te) }\SpecialCharTok{*}\NormalTok{ (a }\SpecialCharTok{+}\NormalTok{ c) }\SpecialCharTok{+}\NormalTok{ b,}
\NormalTok{  otherold }\SpecialCharTok{\textasciitilde{}} \FunctionTok{exp}\NormalTok{(}\SpecialCharTok{{-}}\NormalTok{r }\SpecialCharTok{*}\NormalTok{ tr) }\SpecialCharTok{*}\NormalTok{ d }\SpecialCharTok{*}\NormalTok{ (}\DecValTok{1} \SpecialCharTok{+}\NormalTok{ e }\SpecialCharTok{*}\NormalTok{ te) }\SpecialCharTok{*}\NormalTok{ a }\SpecialCharTok{+}\NormalTok{ b,}
\NormalTok{  npl      }\SpecialCharTok{\textasciitilde{}}\NormalTok{ b,}
\NormalTok{  a }\SpecialCharTok{\textasciitilde{}} \DecValTok{1} \SpecialCharTok{+}\NormalTok{ (}\DecValTok{1} \SpecialCharTok{|}\NormalTok{ participant),}
\NormalTok{  c }\SpecialCharTok{\textasciitilde{}} \DecValTok{1} \SpecialCharTok{+}\NormalTok{ (}\DecValTok{1} \SpecialCharTok{|}\NormalTok{ participant),}
\NormalTok{  d }\SpecialCharTok{\textasciitilde{}} \DecValTok{1} \SpecialCharTok{+}\NormalTok{ (}\DecValTok{1} \SpecialCharTok{|}\NormalTok{ participant),}
\NormalTok{  e }\SpecialCharTok{\textasciitilde{}} \DecValTok{1} \SpecialCharTok{+}\NormalTok{ (}\DecValTok{1} \SpecialCharTok{|}\NormalTok{ participant),}
\NormalTok{  r }\SpecialCharTok{\textasciitilde{}} \DecValTok{1} \SpecialCharTok{+}\NormalTok{ (}\DecValTok{1} \SpecialCharTok{|}\NormalTok{ participant)}
\NormalTok{)}

\NormalTok{fit }\OtherTok{\textless{}{-}} \FunctionTok{bmm}\NormalTok{(}
  \AttributeTok{formula =}\NormalTok{ m3\_pred\_funs,}
  \AttributeTok{model   =}\NormalTok{ m3\_model\_custom,}
  \AttributeTok{data    =}\NormalTok{ sim\_data,}
  \AttributeTok{cores   =} \DecValTok{4}\NormalTok{,}
  \AttributeTok{backend =} \StringTok{"cmdstanr"}\NormalTok{,}
  \AttributeTok{file    =} \FunctionTok{here}\NormalTok{(}\StringTok{"output"}\NormalTok{, }\StringTok{"fit\_m3\_custom\_simple\_N50\_T25"}\NormalTok{)}
\NormalTok{)}
\end{Highlighting}
\end{Shaded}

\section{Recovery Evaluation}\label{recovery-evaluation}

Recovery can be evaluated at two levels. Group-level recovery asks
whether the model recovers the true population means and between-person
variability (i.e., standard deviation of the random intercepts).
Individual-level recovery asks whether each participant's true parameter
value is correctly estimated --- this is critical for researchers who
plan to use \(M^3\) parameter estimates as individual-difference
measures, for example when correlating \(a\) or \(c\) with external
variables.

\subsection{Group-level recovery}\label{group-level-recovery}

We compare the recovered fixed effects to the true generating means. For
each parameter, we extract the posterior median and 95\% credible
interval from \texttt{fixef()} and check two things: (1) bias --- is the
posterior median close to the true value? --- and (2) coverage --- does
the 95\% credible interval contain the true value? We also extract the
between-person SDs from \texttt{brms::VarCorr()} and compare them to the
generating SDs.

\subsection{Individual-level recovery}\label{individual-level-recovery}

Individual-level recovery asks a similar question at the level of
subject-level parameter estimates: for each participant, does the
recovered estimate (fixed + random effect) match their true generating
value? Figure~\ref{fig-supp-indiv-scatter} shows scatter plots of true
versus recovered parameters for a single simulation cell (\(N = 50\), 25
trials per condition). Points close to the diagonal indicate good
recovery; the Pearson correlation (\(r\)) between true and estimated
parameters is annotated in each panel.

\begin{figure}

\caption{\label{fig-supp-indiv-scatter}Individual-level parameter
recovery for the memory updating \(M^3\) (\(N = 50\), 25 trials per
condition). Each panel shows the true generating value (x-axis) versus
the recovered estimate (y-axis) for one parameter. The dashed diagonal
represents perfect recovery; the annotation shows the Pearson
correlation between true and recovered values.}

\centering{

\includegraphics[width=1\linewidth,height=\textheight,keepaspectratio]{../figures/tutorial5_simple_indiv_scatter.pdf}

}

\end{figure}%

Parameters \(a\) and \(c\) show the strongest individual-level recovery
(\(r\) = .84 and .91, respectively). Parameters \(d\) and \(e\) show
moderate recovery (\(r\) = .74 and .70). Recovery is weakest for \(r\)
(\(r\) = .43), which governs the exponential reduction of outdated item
activation over removal time, making it harder to separate from the
other parameters. This is in line with more comprehensive parameter
recovery simulations using the simple choice rule by Göttmann et al.
(\citeproc{ref-goettmann2025recovery}{2025}), who found that \(r\)
requires considerably more data to recover well.

\section{Varying Sample Size and Trial
Count}\label{varying-sample-size-and-trial-count}

The simplified simulation demonstrates the workflow for a single design
cell. To understand how recovery depends on sample size and the number
of trials per condition, the full simulation script
(\texttt{tutorial5\_parameter\_recovery\_updating.R}) varies both
factors across a \(3 \times 3\) grid: \(N \in \{25, 50, 100\}\) and
trials per condition \(\in \{5, 10, 25\}\).

Figure~\ref{fig-supp-group-recovery} shows group-level recovery across
the design grid. Recovered estimates are close to the true values in
most cells, with credible intervals narrowing as \(N\) increases. With
only 5 trials per condition, some parameters (particularly \(e\) and
\(r\)) show noticeable bias, but this improves substantially at 10 and
25 trials.

\begin{figure}

\caption{\label{fig-supp-group-recovery}Group-level parameter recovery
across sample sizes and trial counts. Each panel column represents one
parameter; rows represent different trial counts per condition. Points
and error bars show the recovered posterior median and 95\% credible
interval; dashed horizontal lines mark the true generating values. Black
points indicate that the credible interval covers the true value.}

\centering{

\includegraphics[width=1\linewidth,height=\textheight,keepaspectratio]{../figures/tutorial5_group_recovery.pdf}

}

\end{figure}%

Figure~\ref{fig-supp-heatmap} shows individual-level recovery (Pearson
correlation between true and recovered values) for each parameter across
the same design grid. Individual-level recovery is more demanding than
group-level recovery: \(e\) and \(r\) require both a reasonable sample
size and a sufficient number of trials per condition to reach
correlations above .7.

\begin{figure}

\caption{\label{fig-supp-heatmap}Individual-level recovery heatmap.
Cells show the Pearson correlation between true and recovered
individual-level parameter values for each combination of sample size
(\(N\)) and trials per condition. Darker shading indicates stronger
recovery.}

\centering{

\includegraphics[width=1\linewidth,height=\textheight,keepaspectratio]{../figures/tutorial5_recovery_heatmap.pdf}

}

\end{figure}%

\subsection{Practical guidance}\label{practical-guidance}

These results are based on a single replication per design cell. A
proper simulation study should replicate each cell many times to account
for Monte Carlo variability (\citeproc{ref-talts2018validating}{Talts et
al., 2018}). Recovery quality also depends on the model structure
itself. Parameters that appear in interaction terms or that are only
weakly identified by the design may require more data. As a rule of
thumb, if individual-level correlations fall below .7 for a parameter of
interest, the design likely needs more trials or a larger sample before
that parameter can support individual-differences analyses. Running a
recovery check before collecting empirical data is therefore a practical
necessity for custom models: it reveals whether the planned sample size
and trial count are sufficient, or whether the design needs refinement.

\section*{References}\label{references}
\addcontentsline{toc}{section}{References}

\phantomsection\label{refs}
\begin{CSLReferences}{1}{0}
\bibitem[\citeproctext]{ref-goettmann2025recovery}
Göttmann, J., Frischkorn, G. T., Oberauer, K., Schaefer, S. B., \&
Schubert, A.-L. (2025). \emph{Modeling individual differences in working
memory: {S}ubject-level parameter recovery within the {M}emory
{M}easurement {M}odel framework ({M3})}.
\url{https://doi.org/10.31234/osf.io/945d2}

\bibitem[\citeproctext]{ref-li2025updating}
Li, C., Frischkorn, G. T., \& Oberauer, K. (2025). Updating of
information in working memory: {T}ime course and consequences.
\emph{Cognitive Psychology}, \emph{156}, 101702.
\url{https://doi.org/10.1016/j.cogpsych.2024.101702}

\bibitem[\citeproctext]{ref-oberauer2019simple}
Oberauer, K., \& Lewandowsky, S. (2019). Simple measurement models for
complex working-memory tasks. \emph{Psychological Review},
\emph{126}(6), 880--932. \url{https://doi.org/10.1037/rev0000159}

\bibitem[\citeproctext]{ref-talts2018validating}
Talts, S., Betancourt, M., Simpson, D., Vehtari, A., \& Gelman, A.
(2018). \emph{Validating {B}ayesian inference algorithms with
simulation-based calibration}. \url{https://arxiv.org/abs/1804.06788}

\bibitem[\citeproctext]{ref-wilson2019ten}
Wilson, R. C., \& Collins, A. G. E. (2019). Ten simple rules for the
computational modeling of behavioral data. \emph{eLife}, \emph{8},
e49547. \url{https://doi.org/10.7554/eLife.49547}

\end{CSLReferences}






\end{document}
